
\section*{Open Science}
To support transparency and reproducibility, we have made all source code, simulation tools, and experimental artifacts for \textbf{OptiMAC} available on \textbf{Zenodo}. The repository provides a complete research artifact, including optimization scripts, simulation notebooks, real-world testbed implementations, and recorded results.

\subsubsection*{Contents of the Repository}
\begin{itemize}
    \item \textbf{Optimization Framework}: Python and MATLAB implementations for tag-to-message assignment optimization, including Gurobi-based solvers and parameterized notebooks.
    \item \textbf{Simulation Tools}: MATLAB and Python scripts for analyzing MAC configurations and exploring packet loss scenarios.
    \item \textbf{UDP Demo}: Real-world WiFi and 5G testbed experiments with transmitter/receiver code, and demonstration results.
    \item \textbf{Documentation}: A detailed \texttt{README.md} with setup instructions, usage examples, and a description of the repository structure.
\end{itemize}

\subsubsection*{Accessing the Repository}
The repository is archived on Zenodo and will be accessible under a DOI: \textbf{\href{https://doi.org/10.5281/zenodo.16921174}{doi.org/10.5281/zenodo.16921174}}

During the review process, an anonymized link has been provided to maintain double-blind standards. Upon acceptance, the finalized Zenodo DOI will be published to enable persistent and citable access. 

\subsubsection*{Reproducibility Notes}
\begin{itemize}
    \item \textbf{Environment}: Python 3.9+ is required; MATLAB is optional for running simulation scripts.
    \item \textbf{Dependencies}: Install requirements via \texttt{pip install -r requirements.txt}. Additional libraries such as OpenCV and Gurobi are needed for the UDP demo and optimization experiments.
    \item \textbf{Usage}: Example configurations and step-by-step instructions are provided in the \texttt{README.md}.
\end{itemize}